\chapter{Standards and Data Models}

\section*{Motivation}

Railway alignment data does not exist in isolation. It is embedded in a
landscape of established standards and engineering data models.

Relevant formats include:
\begin{itemize}
\item LandXML,
\item IFC (Industry Foundation Classes),
\item railML,
\item and various proprietary or legacy formats.
\end{itemize}

The AIM framework does not replace these standards. Instead, it provides
a consistent internal representation that can interpret and integrate
them.


\section{General Perspective}

Most existing standards describe railway geometry in a layered or
fragmented way.

\begin{itemize}
\item horizontal alignment,
\item vertical profile,
\item cant (superelevation),
\item additional metadata.
\end{itemize}

In contrast, the AIM framework proposes a unified state-based
representation:
\[
(\gamma(s), \theta(s), \kappa(s), \phi(s)).
\]

This representation serves as a canonical internal model.


\section{LandXML}

\subsection{Overview}

LandXML is one of the most widely used formats for alignment data
exchange.

It typically represents:
\begin{itemize}
\item horizontal geometry (CoordGeom),
\item vertical profiles,
\item cant definitions.
\end{itemize}

\subsection{Characteristics}

\begin{itemize}
\item structured XML format,
\item hierarchical organization,
\item multiple alignments per file,
\item explicit naming of geometric elements.
\end{itemize}

\subsection{Limitations}

Despite its structure, LandXML exhibits several limitations:
\begin{itemize}
\item separation of horizontal and vertical geometry,
\item lack of a unified dynamic interpretation,
\item ambiguity in coordinate reference systems,
\item inconsistent usage across software systems.
\end{itemize}

\subsection{AIM Interpretation}

Within the AIM framework, LandXML is treated as a \emph{fat container}:
\begin{itemize}
\item all available data is parsed,
\item relevant components are extracted,
\item the result is transformed into a unified state model.
\end{itemize}

The conversion step
\[
\text{LandXML} \rightarrow \text{AIM (sparse)}
\]
is essential for further processing.


\section{IFC}

\subsection{Overview}

IFC is a general BIM standard designed for interdisciplinary data
exchange.

In the context of infrastructure, IFC is extended by infrastructure
schemas (IFC Rail, IFC Alignment).

\subsection{Characteristics}

\begin{itemize}
\item object-oriented data model,
\item rich semantic structure,
\item integration of geometry and metadata,
\item support for relationships and hierarchies.
\end{itemize}

\subsection{Challenges}

For alignment modelling, IFC presents several challenges:
\begin{itemize}
\item complex schema structure,
\item multiple representation alternatives,
\item inconsistent implementation across tools,
\item limited support for advanced alignment concepts.
\end{itemize}

\subsection{AIM Interpretation}

The AIM framework interprets IFC data by extracting:
\begin{itemize}
\item alignment geometry,
\item reference systems,
\item semantic relations.
\end{itemize}

AIM acts as a normalization layer that reduces IFC complexity to a
computationally usable form.


\section{railML}

\subsection{Overview}

railML is a domain-specific XML standard for railway systems.

It focuses on:
\begin{itemize}
\item infrastructure,
\item timetable data,
\item rolling stock.
\end{itemize}

\subsection{Relevance for Alignment}

railML provides a more explicit representation of railway topology than
most other formats.

However, detailed geometric modelling is often secondary to network and
operational aspects.

\subsection{AIM Interpretation}

Within AIM, railML is particularly relevant for:
\begin{itemize}
\item network topology,
\item connections between tracks,
\item system-level structure.
\end{itemize}


\section{Legacy and Proprietary Formats}

In practice, a large amount of alignment data exists in proprietary or
legacy formats.

Examples include:
\begin{itemize}
\item TRA / GRA (Technet, Vermessung),
\item spreadsheet-based formats,
\item CAD exports.
\end{itemize}

\subsection{Characteristics}

\begin{itemize}
\item incomplete or implicit structure,
\item missing type information,
\item reliance on conventions rather than formal schema.
\end{itemize}

\subsection{AIM Strategy}

These formats are interpreted using:
\begin{itemize}
\item sniffing and classification,
\item format-specific parsers,
\item reconstruction of implicit structure.
\end{itemize}


\section{Canonical Internal Model}

\begin{definition}[AIM core model]
The AIM framework uses a canonical internal representation based on:
\[
(\gamma(s), \theta(s), \kappa(s), \phi(s)).
\]
\end{definition}

All external formats are mapped to this representation.

\subsection{Advantages}

\begin{itemize}
\item unified handling of geometry and dynamics,
\item compatibility with numerical optimization,
\item independence from specific file formats,
\item extensibility toward new standards.
\end{itemize}


\section{Transformation Pipeline}

The transformation from external data to the AIM model follows a
pipeline structure.

\begin{verbatim}
external file
→ parser
→ fat representation
→ validation
→ sparse conversion
→ AIM state model
\end{verbatim}

This pipeline ensures robustness and traceability of data processing.


\section{Relation to BIM}

\subsection{Conceptual Position}

BIM (Building Information Modelling) aims to integrate geometry,
semantics, and lifecycle information.

However, alignment modelling is often underrepresented or treated as a
secondary aspect.

\subsection{AIM Contribution}

The AIM framework contributes to BIM by:
\begin{itemize}
\item providing a rigorous alignment model,
\item enabling dynamic interpretation,
\item supporting optimization-driven design.
\end{itemize}

In this sense, AIM can be interpreted as a specialization of BIM for
alignment-centric infrastructure modelling.


\section{Interoperability}

A key requirement for practical application is interoperability with
existing tools and workflows.

The AIM framework supports interoperability through:
\begin{itemize}
\item import adapters,
\item export converters,
\item intermediate canonical representation.
\end{itemize}


\section{Discussion}

Existing standards focus primarily on data exchange and documentation.

The AIM framework complements these standards by introducing:
\begin{itemize}
\item a computationally consistent model,
\item a dynamic interpretation,
\item integration with optimization methods.
\end{itemize}

Thus, AIM bridges the gap between static data representation and active
engineering design.


\section{Summary}

\begin{summary}
Railway alignment data is governed by a variety of standards, each with
its own strengths and limitations.

The AIM framework provides a unified internal representation that allows
these heterogeneous sources to be interpreted consistently.

By combining standard-based data exchange with a state-based,
optimization-ready model, AIM enables a transition from static data
handling to dynamic and computational alignment design.
\end{summary}
